\chapter{Клеточные автоматы, обратимость и представление Маделунга}
\label{app:background}

\section{Зачем это приложение}

Структурное определение клеточного автомата, дискретной фазы и координат $(dh,dt)$
дано в конце спуска (гл.~\ref{ch:descent}, \cref{sec:descent-ca,sec:coords-dh-dt}).
Здесь --- исторический и учебный контекст: классические имена, Тьюринг-полнота,
развёрнутый Маделунг и литература.
Читатель, идущий по основному тексту, может пропустить приложение.

\section{Исторические клеточные автоматы}

Идея восходит к Уламу и фон Нейману (1940-е; см.\ \cite{vonneumann1966}):
сложное поведение из локальных правил на решётке.
«Игра жизни'' Конвея \cite{gardner1970} и Rule~110 \cite{cook2004}
показывают, что даже простые булевы правила Тьюринг-полны.
Наша модель~$M$ --- \emph{другой} элемент того же класса:
обратимый спинорный автомат с регистром $\spaceZ_{\phaseRingCount}[i]$, не подобранный для картинки.

\section{Тьюринг-полнота, дифференциально-интегральные операторы и физический тезис Чёрча--Тьюринга}
\label{sec:app-universality}

\emph{Тьюринг-полнота} клеточного автомата означает его способность имитировать произвольную машину Тьюринга при достаточном пространственном объёме (классические конструкции в «Игре жизни» Конвея \cite{gardner1970}, правиле~110 \cite{cook2004} и клеточном автомате «муравей Лэнгтона» \cite{langton1986}, где одиночный локальный агент на двумерной решётке реализует универсальное вычисление).
Однако при сопоставлении клеточных автоматов с уравнениями математической физики долгое время господствовал обратный взгляд, восходящий к работе Омохундро \cite{omohundro1984}:
непрерывные уравнения в частных производных рассматривались как первичная реальность, а дискретный автомат --- лишь как грубая конечная аппроксимация континуума.

В настоящей монографии направление этой стрелки обращено:
\begin{enumerate}[leftmargin=*]
\item \textbf{Единое порождение дифференциальных и интегральных операторов:}
  одношаговая разность по ближайшим соседям звезды~$N_{12}$ задаёт дискретные дифференциальные операторы (градиент, ротор, дивергенцию и лапласиан), тогда как $R$-кратная итерация локального шага по путям светового конуса за $R$ тактов автоматически вычисляет многомерную свёртку с функцией Грина $\Phi = G_{\mathrm{FCC}}^{(2R)} * z$ (\cref{ch:macro}), то есть решает интегральные уравнения поля без нелокального интегрирования.
\item \textbf{Физическое уточнение тезиса Чёрча--Тьюринга:}
  с одной стороны, предел Бекенштейна $\caBitBudget = 2\pi/\ln 2$ исключает существование континуального «гипервычисления» на вещественных числах в конечном объёме;
  с другой стороны, последовательная одноленточная машина Тьюринга со стиранием символов не является адекватной моделью физического времени.
  Фундаментальным вычислителем природы выступает однородный обратимый трёхмерный клеточный автомат на решётке кубооктаэдра~\latGCKshort\ над конечным кольцом $\spaceZ_{\phaseRingCount}[i]$, выполняющий синхронный шаг на всех узлах за $O(1)$ операций в такт без энтропийного стирания информации ($g^{-1}\circ g = \mathrm{id}$).
\end{enumerate}

\begin{remark}[Три уровня описания]
\begin{enumerate}
\item \textbf{Класс клеточных автоматов} --- универсальность по Тьюрингу при строгой локальности и дискретности.
\item \textbf{Закон эволюции~$g$} --- единственный обратимый элемент этого класса на решётке~\latGCK, удовлетворяющий абсолютным условиям $A_1$--$A_{16}$.
\item \textbf{Макроуровень~$T$} --- крупномасштабное биномиальное осреднение~$\mathcal{B}$, порождающее гладкие дифференциальные и интегральные уравнения поля и статистику Борна \eqref{eq:born}.
\end{enumerate}
\end{remark}

\section{Обратимость: Тоффоли и Фредкин}

На~$M$ требуется $g^{-1}(g(\symCapitalPsi))=\symCapitalPsi$ --- иначе унитарность \eqref{eq:norm} нарушается на шаге.
Фредкин \cite{fredkin1990} и Тоффоли с Марголусом \cite{toffoli1987} развивали обратимые CA
как модель вычисления без энтропийного стирания бита (см.\ также \cite{bennett1973,landauer1961}).
Физическая мотивация в монографии та же, что в спуске:
микрошаг --- перестановка конечного множества состояний; необратимость измерения --- на~$T$.

\section{Представление Маделунга}

В 1927 году Маделунг \cite{madelung1927} переписал уравнение Шрёдингера в гидродинамической форме:
$\wavePsi=\sqrt{\massDensity}\,e^{iS/\planckHbar}$, откуда неразрывность $\opPartialT\chargeRho+\opNabla\cdot(\symRho\fieldV)=0$
и квантовое давление $Q\propto\opNabla^2\sqrt{\massDensity}/\sqrt{\massDensity}$.

Маделунг --- язык \emph{непрерывного}~$T$.
На~$M$ нет гладкого $\opPartialT\chargeRho+\opNabla\cdot\fieldJ=0$ как фундаментального закона;
норма сохраняется дискретно, а гидродинамика появляется после осреднения
(гл.~\ref{ch:macro}, \ref{ch:matter}).
Дискретизация фазы в спуске (\cref{sec:descent-phase-disc}) --- не «доказательство Маделунгом'',
а замена $\phaseVarphi\in\spaceR$ конечным $\phaseVarphi_{\mathrm{disc}}$ при том же смысле $|z|^2$ и фазы.

\begin{remark}[Частые заблуждения]
\begin{enumerate}
\item \textbf{«Маделунг доказывает дискретность.''} Нет --- он переписывает непрерывное уравнение.
\item \textbf{«КА = Игра жизни.''} Игра жизни --- частный необратимый пример.
\item \textbf{«На $M$ --- неразрывность Маделунга.''} Гладкость --- свойство~$T$ при $\lengthScale\gg\caStepSpace$.
\end{enumerate}
\end{remark}

\section{Сводная таблица}

\begin{table}[ht]
\centering
\caption{От классических понятий к главам монографии}
\label{tab:app-map}
\small
\begin{tabular}{@{}p{3.2cm}p{4.2cm}p{5.8cm}@{}}
\toprule
Понятие & Классика & У нас \\
\midrule
Клеточный автомат & Улама, von Neumann & Спуск \cref{sec:descent-ca}; закон~$g$ \\
Тьюринг-полнота & Конвей, Rule~110 & Класс CA; \S\ref{sec:app-universality} \\
Обратимость & Фредкин, Тоффоли & $g^{-1}$, гл.~\ref{ch:evolution} \\
Дискретная фаза & --- & Спуск \cref{sec:descent-phase-disc} \\
$(dh,dt)$ & --- & Спуск \cref{sec:coords-dh-dt} \\
Скорость света CA & $\caSignalSpeed=1$ cell/tick & $\caSignalSpeed=\caStepSpace/\caStepTime$; см. \cref{sec:ca-terminology,sec:two-c} \\
Маделунг & 1927, $\symRho$, $\fieldJ$ & Только~$T$ \\
\bottomrule
\end{tabular}
\end{table}

\section{Литература}

Развёрнутый библиографический список по темам CA, обратимости и цифровой физики ---
в гл.\ \ref{ch:literature}.
Обзорный том Вольфрама \cite{wolfram2002} удобен для ориентации в классе правил;
конкретный закон~$g$ монографии его классификации не следует.
